\documentclass[10pt,a4paper]{article}
\usepackage[T1]{fontenc}
\usepackage[utf8]{inputenc}
\usepackage[czech]{babel}
\usepackage{lmodern}
\usepackage{microtype}
\usepackage{geometry}
\geometry{left=14mm,right=14mm,top=18mm,bottom=18mm}
\usepackage[table]{xcolor}
\usepackage{array,longtable,tabularx,booktabs}
\usepackage{fancyhdr,lastpage}
\usepackage{ragged2e}
\usepackage{hyperref}
\hypersetup{hidelinks,pdftitle={Rozpis výuky matematiky – 452 hodin podle tematických plánů a skutečné praxe}}
\setlength{\parindent}{0pt}
\setlength{\parskip}{5pt}
\setlength{\tabcolsep}{4pt}
\renewcommand{\arraystretch}{1.08}
\setlength{\headheight}{18pt}
\pagestyle{fancy}\fancyhf{}
\fancyhead[L]{\footnotesize Matematika -- rozpis výuky}
\fancyhead[R]{\footnotesize 1.--4. ročník SPŠ}
\fancyfoot[L]{\footnotesize Rozpis hodin podle tematického plánu}
\fancyfoot[R]{\footnotesize \thepage/\pageref{LastPage}}
\renewcommand{\headrulewidth}{0.35pt}
\newcommand{\testr}[1]{\textbf{#1}}
\begin{document}
\begin{center}
{\LARGE\bfseries Rozpis výuky matematiky}\\[3mm]
{\large 1.--4. ročník střední průmyslové školy}\\[3mm]
{\small Přehled plánované výuky podle tematických plánů SP\v{S} a reálných zápisů výuky}
\end{center}
\vspace{3mm}
\begin{tabularx}{\textwidth}{@{}Xrrrr@{}}\toprule
\textbf{Ročník} & \textbf{Hodin} & \textbf{Krátké testy} & \textbf{Souhrnné} & \textbf{Testů celkem}\\\midrule
1. ročník & 136 & 8 & 4 & 12\\
2. ročník & 102 & 8 & 4 & 12\\
3. ročník & 102 & 8 & 4 & 12\\
4. ročník & 112 & 8 & 4 & 12\\\midrule
\textbf{Celkem} & \textbf{452} & \textbf{32} & \textbf{16} & \textbf{48}\\\bottomrule
\end{tabularx}
\vspace{4mm}
\section*{Charakter dokumentu}
\begin{itemize}\setlength{\itemsep}{2pt}
\item Dokument obsahuje rozpis výuky matematiky pro 1.--4. ročník SP\v{S} v celkové dotaci \textbf{452 hodin}.
\item Rozpis je členěn po ročnících a dále po měsících. Každý očíslovaný řádek představuje \textbf{jednu plánovanou vyučovací hodinu}.
\item Zkratka \textbf{TP} odkazuje na odpovídající tematický okruh tematického plánu. U testů je označení uvedeno přímo v náplni hodiny.
\item V každém čtvrtletí jsou zařazeny \textbf{tři testy}: dva průběžné testy (zpravidla 15--20 minut) a jedna souhrnná písemná práce.
\end{itemize}

\section*{Z čeho rozpis vychází}
\begin{itemize}\setlength{\itemsep}{2pt}
\item Hodinová dotace vychází z tematických plánů SP\v{S}: \textbf{136 + 102 + 102 + 112 hodin}.
\item Obsah a pořadí témat zohledňují reálné zápisy výuky v dostupných třídách, zejména v 1.A (2024/25), 2.A (2025/26), 3.A (2025/26) a 3.D (2025/26).
\item Historické zápisy slouží jako podklad pro reálné pořadí témat a odhad časové náročnosti; nejde o doslovný přepis třídních knih.
\end{itemize}

\section*{Podkladové dokumenty}
\begin{itemize}\setlength{\itemsep}{2pt}
\item tematické plány SP\v{S},
\item zápisy výuky.
\end{itemize}

\section*{Odlišnosti oproti tematickým plánům}
\begin{itemize}\setlength{\itemsep}{2pt}
\item \textbf{1. ročník:} výrazněji je rozpracován blok výrokové logiky a množin; pořadí témat odpovídá reálné výuce a zachovává již vytvořené první přípravy.
\item \textbf{2. ročník:} více prostoru dostávají exponenciální rovnice, substituce a procvičování navazujících algebraických úprav.
\item \textbf{3. ročník:} pořadí tematických celků vychází ze skutečné výuky ve třídách 3.A a 3.D; zohledněno je i odlišné tempo obou tříd.
\item \textbf{4. ročník:} protože nejsou k dispozici úplné zápisy výuky, vychází rozpis především z tematického plánu a z již dříve připravené osnovy.
\end{itemize}
\clearpage

\section*{1. ročník – logika, množiny, algebra a funkce}
\textit{Dotace 136 hodin. Čtvrtletní celky a tři testy v každém čtvrtletí jsou součástí hodinové dotace.}
\vspace{2mm}
\footnotesize
\begin{longtable}{@{}>{\RaggedLeft\arraybackslash}p{9mm}>{\RaggedRight\arraybackslash}p{15mm}>{\RaggedRight\arraybackslash}p{150mm}@{}}
\toprule \textbf{Hod.} & \textbf{TP} & \textbf{Náplň vyučovací hodiny} \\\midrule\endfirsthead
\toprule \textbf{Hod.} & \textbf{TP} & \textbf{Náplň vyučovací hodiny} \\\midrule\endhead
\bottomrule\endfoot
\addlinespace[1mm]\multicolumn{3}{@{}l}{\footnotesize\textbf{1. čtvrtletí} -- 2 průběžné testy + 1 souhrnná práce}\\[0.5mm]
\addlinespace[1.1mm]\multicolumn{2}{@{}l}{\textbf{Září}} & \multicolumn{1}{r@{}}{\textbf{12 hodin}}\\
\midrule
1 & 1.1 & Úvodní hodina: pravidla práce, hodnocení, pomůcky a organizace výuky. \\
2 & 1.1 & Co je výrok? Výroky a nevýroky, pravdivostní hodnota, jednoduchý a složený výrok. \\
3 & 1.1 & Pravdivost složených výroků: negace, konjunkce, disjunkce, implikace, ekvivalence; tabulky. \\
4 & 1.1 & Obecný a existenční kvantifikátor, otevřená formule, negace kvantifikovaných výroků. \\
5 & 1.2 & Číselné obory: přirozená a celá čísla, znaménka, přednost operací. \\
6 & 1.2 & Racionální čísla: zlomky, rozšiřování, krácení a porovnávání. \\
7 & 1.2 & Reálná čísla: iracionální čísla, číselná osa, zaokrouhlování. \\
8 & 1.2 & Intervaly: krajní body, otevřený a uzavřený interval. \\
9 & 1.2 & Absolutní hodnota čísla jako vzdálenost na ose; jednoduché úlohy. \\
10 & T1 & \testr{Průběžný test 1 (15–20 min): Výroková logika a číselné obory. Zbytek hodiny: krátký rozbor a navazující procvičování.} \\
11 & 1.2 & Číselné obory a výroky: pravdivost podle zvolené množiny; oprava chyb z testu. \\
12 & 1.1–1.2 & Množina a její prvky: určení výčtem a vlastností; prázdná množina. \\
\addlinespace[1.1mm]\multicolumn{2}{@{}l}{\textbf{Říjen}} & \multicolumn{1}{r@{}}{\textbf{16 hodin}}\\
\midrule
13 & 1.1–1.2 & Vztahy mezi množinami: podmnožina a rovnost množin. \\
14 & 1.1–1.2 & Vennovy diagramy: znázornění dvou a tří množin. \\
15 & 1.1–1.2 & Sjednocení a průnik množin: význam a zápis. \\
16 & 1.1–1.2 & Rozdíl a doplněk množiny vzhledem k základní množině. \\
17 & 1.1–1.2 & Množinové operace a jejich vlastnosti; kontrola Vennovým diagramem. \\
18 & 1.1–1.2 & Číselné množiny a intervaly: sjednocení a průnik intervalů. \\
19 & 1.1–1.2 & Slovní úlohy na množiny: třídění prvků a počítání počtů. \\
20 & 1.1–1.2 & Opakování množin: zápisy, diagramy, kombinované úlohy. \\
21 & T2 & \testr{Průběžný test 2 (15–20 min): Množiny, intervaly a množinové operace. Zbytek hodiny: krátký rozbor a navazující procvičování.} \\
22 & 1.3 & Pravoúhlý trojúhelník: Pythagorova věta, výpočty délek a úhlů. \\
23 & 1.3 & Sinus, kosinus a tangens ostrého úhlu; poměry stran v pravoúhlém trojúhelníku. \\
24 & 1.3 & Výpočty výšek a vzdáleností pomocí goniometrie; aplikace a kontrola jednotek. \\
25 & 2.1–2.4 & Mocnina s přirozeným exponentem: význam, znaménko, závorky. \\
26 & 2.1–2.4 & Pravidla pro součin a podíl mocnin stejného základu. \\
27 & 2.1–2.4 & Mocnina mocniny, součinu a podílu; odhalování chybných pravidel. \\
28 & 2.1–2.4 & Nulový a záporný celočíselný exponent; převrácená hodnota. \\
\addlinespace[1.1mm]\multicolumn{2}{@{}l}{\textbf{Listopad}} & \multicolumn{1}{r@{}}{\textbf{6 hodin}}\\
\midrule
29 & 2.1–2.4 & Vědecký zápis čísel a řádové odhady. \\
30 & 2.1–2.4 & Druhá a třetí odmocnina; vztah k druhé a třetí mocnině. \\
31 & 2.1–2.4 & Mocniny s racionálním exponentem; převody na odmocniny. \\
32 & P1 & \testr{Čtvrtletní písemná práce: Výroková logika, množiny, goniometrie a mocniny.} \\
33 & 2.1–2.4 & Základní věty pro odmocniny včetně podmínek platnosti. \\
34 & 2.1–2.4 & Vynášení a částečné odmocňování, slučování odmocnin. \\
\addlinespace[1mm]\multicolumn{3}{@{}l}{\footnotesize\textbf{2. čtvrtletí} -- 2 průběžné testy + 1 souhrnná práce}\\[0.5mm]
35 & 2.1–2.4 & Usměrňování jmenovatele s jedním odmocninovým členem. \\
36 & 2.1–2.4 & Usměrňování jmenovatele se sdruženým výrazem. \\
37 & 2.1–2.4 & Složené výrazy s mocninami a odmocninami – postup a podmínky. \\
38 & 2.1–2.4 & Slovní a technické úlohy s mocninami a odmocninami. \\
39 & 2.1–2.4 & Opakování mocnin a odmocnin; cílené procvičení před testem. \\
40 & T3 & \testr{Průběžný test 3 (15–20 min): Mocniny a odmocniny. Zbytek hodiny: krátký rozbor a navazující procvičování.} \\
41 & 3.1–3.5 & Algebraický výraz, proměnná, koeficient, člen a stupeň mnohočlenu. \\
42 & 3.1–3.5 & Sčítání a odčítání mnohočlenů; roznásobování závorek. \\
43 & 3.1–3.5 & Násobení jednočlenu mnohočlenem; znaménka a pořadí úprav. \\
44 & 3.1–3.5 & Násobení dvou mnohočlenů; uspořádání výsledku. \\
\addlinespace[1.1mm]\multicolumn{2}{@{}l}{\textbf{Prosinec}} & \multicolumn{1}{r@{}}{\textbf{12 hodin}}\\
\midrule
45 & 3.1–3.5 & Dělení mnohočlenu jednočlenem; opakování násobení. \\
46 & 3.1–3.5 & Dělení mnohočlenů mnohočlenem: princip algoritmu. \\
47 & 3.1–3.5 & Dělení mnohočlenů se zbytkem; kontrola násobením. \\
48 & 3.1–3.5 & Rozklad mnohočlenu vytýkáním společného činitele. \\
49 & 3.1–3.5 & Vzorce pro druhou mocninu dvojčlenu a rozdíl čtverců. \\
50 & 3.1–3.5 & Kombinace vytýkání a vzorců při rozkladu. \\
51 & 3.1–3.5 & Součet a rozdíl třetích mocnin; vhodná volba vzorce. \\
52 & 3.1–3.5 & Lomené algebraické výrazy: podmínky, krácení. \\
53 & 3.1–3.5 & Sčítání a odčítání lomených algebraických výrazů. \\
54 & 3.1–3.5 & Násobení a dělení lomených výrazů; součinové rozklady. \\
55 & 3.1–3.5 & Složené lomené výrazy; společné podmínky. \\
56 & 3.1–3.5 & Smíšené úpravy polynomů a lomených výrazů. \\
\addlinespace[1.1mm]\multicolumn{2}{@{}l}{\textbf{Leden}} & \multicolumn{1}{r@{}}{\textbf{16 hodin}}\\
\midrule
57 & T4 & \testr{Průběžný test 4 (15–20 min): Algebraické výrazy a rozklady. Zbytek hodiny: krátký rozbor a navazující procvičování.} \\
58 & 3.1–3.5 & Opakování algebraických výrazů; důkaz správnosti dosazením. \\
59 & 4.1–4.4 & Funkce jako přiřazení: argument a funkční hodnota, předpis. \\
60 & 4.1–4.4 & Definiční obor a obor hodnot z předpisu a grafu. \\
61 & 4.1–4.4 & Graf funkce: intervaly růstu a poklesu, průsečíky s osami. \\
62 & 4.1–4.4 & Základní vlastnosti funkcí: monotonie, sudost, lichost. \\
63 & 4.1–4.4 & Lineární funkce a její předpis; graf jako přímka. \\
64 & 4.1–4.4 & Směrnice a úsek na ose: geometrický význam parametrů. \\
65 & 4.1–4.4 & Sestrojení grafu lineární funkce z předpisu a tabulky. \\
66 & 4.1–4.4 & Konstantní funkce; vztah k lineární funkci. \\
67 & 4.1–4.4 & Lineární funkce v technických a ekonomických úlohách. \\
68 & 4.1–4.4 & Grafické porovnávání dvou lineárních funkcí. \\
69 & P2 & \testr{Pololetní písemná práce: Algebraické výrazy a základy funkcí.} \\
70 & 4.1–4.4 & Funkce s absolutní hodnotou: graf y=|x| a základní transformace. \\
71 & 4.1–4.4 & Samostatné procvičování grafů a vlastností lineárních funkcí. \\
72 & 4.1–4.4 & Opakování funkcí; přechod od grafu k rovnici. \\
\addlinespace[1mm]\multicolumn{3}{@{}l}{\footnotesize\textbf{3. čtvrtletí} -- 2 průběžné testy + 1 souhrnná práce}\\[0.5mm]
\addlinespace[1.1mm]\multicolumn{2}{@{}l}{\textbf{Únor}} & \multicolumn{1}{r@{}}{\textbf{12 hodin}}\\
\midrule
73 & 4.5–4.11 & Lineární rovnice: ekvivalentní úpravy a zkouška. \\
74 & 4.5–4.11 & Rovnice se závorkami a zlomky; postupné úpravy. \\
75 & 4.5–4.11 & Lineární nerovnice: změna znaménka při násobení záporným číslem. \\
76 & 4.5–4.11 & Rovnice s parametrem: jeden kořen, žádný kořen, nekonečně mnoho. \\
77 & 4.5–4.11 & Soustava dvou lineárních rovnic: dosazovací metoda. \\
78 & 4.5–4.11 & Soustava dvou lineárních rovnic: sčítací metoda. \\
79 & 4.5–4.11 & Soustavy rovnic: geometrická interpretace a počet řešení. \\
80 & 4.5–4.11 & Rovnice s absolutní hodnotou: vzdálenost na číselné ose. \\
81 & 4.5–4.11 & Rovnice s absolutní hodnotou: rozdělení na případy. \\
82 & 4.5–4.11 & Soustavy lineárních rovnic se třemi neznámými. \\
83 & 4.5–4.11 & Soustavy rovnic: slovní úlohy a kontrola dosazením. \\
84 & T5 & \testr{Průběžný test 5 (15–20 min): Lineární rovnice a nerovnice. Zbytek hodiny: krátký rozbor a navazující procvičování.} \\
\addlinespace[1.1mm]\multicolumn{2}{@{}l}{\textbf{Březen}} & \multicolumn{1}{r@{}}{\textbf{16 hodin}}\\
\midrule
85 & 4.5–4.11 & Smíšené lineární rovnice s parametrem: rozbor případů. \\
86 & 4.5–4.11 & Lineární nerovnice se zlomky a závorkami. \\
87 & 4.5–4.11 & Soustavy lineárních nerovnic: průnik množin řešení. \\
88 & 4.5–4.11 & Intervalový zápis řešení nerovnic a soustav. \\
89 & 4.5–4.11 & Nerovnice s absolutní hodnotou: geometrický význam. \\
90 & 4.5–4.11 & Nerovnice s absolutní hodnotou: rozklad podle znaménka. \\
91 & 4.5–4.11 & Rovnice se dvěma absolutními hodnotami: kritické body. \\
92 & 4.5–4.11 & Nerovnice se dvěma absolutními hodnotami: jednoduché úlohy. \\
93 & 4.5–4.11 & Lineární rovnice a nerovnice: kombinované procvičení. \\
94 & 4.5–4.11 & Slovní úlohy vedoucí na lineární rovnice. \\
95 & 4.5–4.11 & Slovní úlohy vedoucí na soustavy rovnic. \\
96 & 4.5–4.11 & Vyjádření neznámé z technického vzorce a návaznost na rovnice. \\
97 & 4.5–4.11 & Vyjádření neznámé z lomeného technického vzorce; podmínky. \\
98 & T6 & \testr{Průběžný test 6 (15–20 min): Soustavy a rovnice s absolutní hodnotou. Zbytek hodiny: krátký rozbor a navazující procvičování.} \\
99 & 4.5–4.11 & Procvičení soustav rovnic a nerovnic – volba metody. \\
100 & 4.5–4.11 & Opakování lineárních rovnic s absolutní hodnotou. \\
\addlinespace[1.1mm]\multicolumn{2}{@{}l}{\textbf{Duben}} & \multicolumn{1}{r@{}}{\textbf{12 hodin}}\\
\midrule
101 & 4.5–4.11 & Diagnostika typických chyb a oprava předchozích úloh. \\
102 & 4.5–4.11 & Souhrnné opakování lineárních rovnic, nerovnic a soustav. \\
103 & 5.1–5.9 & Kvadratická funkce: parabola, definiční obor a obor hodnot. \\
104 & 5.1–5.9 & Graf kvadratické funkce y=ax²; vliv parametru a. \\
105 & 5.1–5.9 & Posuny paraboly a vrcholový tvar; souvislost s osou souměrnosti. \\
106 & 5.1–5.9 & Doplnění na čtverec: získání vrcholového tvaru. \\
107 & 5.1–5.9 & Kvadratický trojčlen a průsečíky paraboly s osami. \\
108 & 5.1–5.9 & Kvadratická funkce v obecném tvaru: konstrukce grafu. \\
109 & 5.1–5.9 & Grafy kvadratických funkcí – samostatná práce. \\
110 & P3 & \testr{Čtvrtletní písemná práce: Lineární funkce a rovnice, úvod kvadratických funkcí.} \\
111 & 5.1–5.9 & Neúplné kvadratické rovnice: vytýkání a odmocnění. \\
112 & 5.1–5.9 & Úplná kvadratická rovnice: diskriminant a vzorec. \\
\addlinespace[1mm]\multicolumn{3}{@{}l}{\footnotesize\textbf{4. čtvrtletí} -- 2 průběžné testy + 1 souhrnná práce}\\[0.5mm]
\addlinespace[1.1mm]\multicolumn{2}{@{}l}{\textbf{Květen}} & \multicolumn{1}{r@{}}{\textbf{12 hodin}}\\
\midrule
113 & 5.1–5.9 & Kvadratická rovnice: rozklad na součin a nulový součin. \\
114 & 5.1–5.9 & Viètovy vztahy: součet a součin kořenů. \\
115 & 5.1–5.9 & Sestavení kvadratické rovnice z daných kořenů. \\
116 & 5.1–5.9 & Grafické a algebraické řešení kvadratických rovnic – srovnání. \\
117 & 5.1–5.9 & Rovnice vedoucí na kvadratickou rovnici: jednoduchá substituce. \\
118 & 5.1–5.9 & Soustava lineární a kvadratické rovnice o dvou neznámých. \\
119 & 5.1–5.9 & Soustavy rovnic s parabolou: geometrický význam průsečíků. \\
120 & 5.1–5.9 & Rovnice s odmocninou: podmínky a ověřování kořenů. \\
121 & T7 & \testr{Průběžný test 7 (15–20 min): Kvadratické rovnice a grafy. Zbytek hodiny: krátký rozbor a navazující procvičování.} \\
122 & 5.1–5.9 & Rovnice se dvěma odmocninami: izolace a kontrola. \\
123 & 5.1–5.9 & Kvadratická rovnice s parametrem: zvláštní případy. \\
124 & 5.1–5.9 & Kvadratická nerovnice: odečítání řešení z grafu paraboly. \\
\addlinespace[1.1mm]\multicolumn{2}{@{}l}{\textbf{Červen}} & \multicolumn{1}{r@{}}{\textbf{12 hodin}}\\
\midrule
125 & 5.1–5.9 & Kvadratická nerovnice: metoda nulových bodů a znamének. \\
126 & 5.1–5.9 & Nerovnice s kvadratickým trojčlenem a intervaly. \\
127 & 5.1–5.9 & Souhrnné úlohy na kvadratické rovnice a nerovnice. \\
128 & 5.1–5.9 & Slovní úlohy vedoucí na kvadratickou rovnici. \\
129 & 5.1–5.9 & Grafická optimalizace jednoduché kvadratické funkce. \\
130 & T8 & \testr{Průběžný test 8 (15–20 min): Kvadratické nerovnice a odmocninové rovnice. Zbytek hodiny: krátký rozbor a navazující procvičování.} \\
131 & 5.1–5.9 & Technické aplikace: obsah a rozměry, kontrola fyzikální smysluplnosti. \\
132 & 5.1–5.9 & Opakování kvadratických funkcí: graf, vrchol, nulové body. \\
133 & 5.1–5.9 & Opakování kvadratických rovnic: metody a volba postupu. \\
134 & 5.1–5.9 & Oprava typických chyb z kvadratických nerovnic. \\
135 & P4 & \testr{Pololetní písemná práce: Souhrnné učivo 1. ročníku.} \\
136 & 5.1–5.9 & Závěrečné procvičení a shrnutí prvního ročníku. \\
\end{longtable}
\normalsize
\clearpage
\section*{2. ročník – exponenciální a logaritmické funkce, goniometrie, geometrie}
\textit{Dotace 102 hodin. Čtvrtletní celky a tři testy v každém čtvrtletí jsou součástí hodinové dotace.}
\vspace{2mm}
\footnotesize
\begin{longtable}{@{}>{\RaggedLeft\arraybackslash}p{9mm}>{\RaggedRight\arraybackslash}p{15mm}>{\RaggedRight\arraybackslash}p{150mm}@{}}
\toprule \textbf{Hod.} & \textbf{TP} & \textbf{Náplň vyučovací hodiny} \\\midrule\endfirsthead
\toprule \textbf{Hod.} & \textbf{TP} & \textbf{Náplň vyučovací hodiny} \\\midrule\endhead
\bottomrule\endfoot
\addlinespace[1mm]\multicolumn{3}{@{}l}{\footnotesize\textbf{1. čtvrtletí} -- 2 průběžné testy + 1 souhrnná práce}\\[0.5mm]
\addlinespace[1.1mm]\multicolumn{2}{@{}l}{\textbf{Září}} & \multicolumn{1}{r@{}}{\textbf{9 hodin}}\\
\midrule
1 & 1.1–1.4 & Úvod do druhého ročníku: pravidla, opakování základních funkcí. \\
2 & 1.1–1.4 & Funkce: definiční obor, obor hodnot a graf z prvního ročníku. \\
3 & 1.1–1.4 & Vlastnosti funkcí: monotonie, prostota a intervaly. \\
4 & 1.1–1.4 & Grafické transformace funkcí: posun a změna měřítka. \\
5 & 1.1–1.4 & Krátké opakování funkcí: samostatná práce s grafem. \\
6 & 1.1–1.4 & Prostá funkce a podmínky existence inverzní funkce. \\
7 & 1.1–1.4 & Inverzní funkce: výměna argumentu a hodnoty; symetrie grafu. \\
8 & 1.1–1.4 & Určení předpisu inverzní funkce a kontrola oborů. \\
9 & T1 & \testr{Průběžný test 1 (15–20 min): Funkce a inverzní funkce. Zbytek hodiny: krátký rozbor a navazující procvičování.} \\
\addlinespace[1.1mm]\multicolumn{2}{@{}l}{\textbf{Říjen}} & \multicolumn{1}{r@{}}{\textbf{12 hodin}}\\
\midrule
10 & 1.1–1.4 & Mocninné funkce: kladné a záporné celočíselné exponenty. \\
11 & 1.5, 1.7–1.8 & Exponenciální funkce: definice, základ, graf a vlastnosti. \\
12 & 1.5, 1.7–1.8 & Srovnání růstu a poklesu; základy větší než jedna a mezi nulou a jednou. \\
13 & 1.5, 1.7–1.8 & Exponenciální rovnice: shodné základy a porovnání exponentů. \\
14 & 1.5, 1.7–1.8 & Úpravy exponentů; převod mocnin na společný základ. \\
15 & 1.5, 1.7–1.8 & Exponenciální rovnice pomocí vytýkání mocniny. \\
16 & 1.5, 1.7–1.8 & Exponenciální rovnice se substitucí t=a\textasciicircum{}x>0: první příklady. \\
17 & 1.5, 1.7–1.8 & Kvadratická rovnice po substituci: návrat k x. \\
18 & 1.5, 1.7–1.8 & Substituce ve složitějších exponenciálních rovnicích. \\
19 & 1.5, 1.7–1.8 & Grafy exponenciálních funkcí: konstrukce a čtení. \\
20 & T2 & \testr{Průběžný test 2 (15–20 min): Exponenciální funkce a grafy. Zbytek hodiny: krátký rozbor a navazující procvičování.} \\
21 & 1.5, 1.7–1.8 & Exponenciální rovnice se součty mocnin; účelná úprava. \\
\addlinespace[1.1mm]\multicolumn{2}{@{}l}{\textbf{Listopad}} & \multicolumn{1}{r@{}}{\textbf{9 hodin}}\\
\midrule
22 & 1.5, 1.7–1.8 & Soustavy exponenciálních rovnic: základní postup. \\
23 & 1.5, 1.7–1.8 & Soustavy exponenciálních rovnic se společným základem. \\
24 & 1.5, 1.7–1.8 & Exponenciální rovnice: procvičování vytýkání a substituce. \\
25 & 1.5, 1.7–1.8 & Exponenciální nerovnice: monotonie podle základu. \\
26 & 1.5, 1.7–1.8 & Exponenciální nerovnice s úpravou společného základu. \\
27 & 1.5, 1.7–1.8 & Soustavy exponenciálních rovnic: opakování a rozbor. \\
28 & 1.5, 1.7–1.8 & Technické aplikace exponenciálního růstu a poklesu. \\
29 & P1 & \testr{Čtvrtletní písemná práce: Funkce a exponenciální rovnice.} \\
30 & 1.5, 1.7–1.8 & Exponenciální rovnice: smíšené příklady před souhrnným testem. \\
\addlinespace[1mm]\multicolumn{3}{@{}l}{\footnotesize\textbf{2. čtvrtletí} -- 2 průběžné testy + 1 souhrnná práce}\\[0.5mm]
31 & 1.5, 1.7–1.8 & Soustavy exponenciálních rovnic: rozšířené procvičení. \\
32 & 1.6, 1.9–1.10 & Logaritmus jako inverzní operace k umocňování; podmínky. \\
33 & 1.6, 1.9–1.10 & Dekadický a přirozený logaritmus; základní hodnoty. \\
\addlinespace[1.1mm]\multicolumn{2}{@{}l}{\textbf{Prosinec}} & \multicolumn{1}{r@{}}{\textbf{9 hodin}}\\
\midrule
34 & 1.6, 1.9–1.10 & Věty o logaritmu součinu, podílu a mocniny. \\
35 & 1.6, 1.9–1.10 & Logaritmická funkce: graf, definiční obor a monotonie. \\
36 & 1.6, 1.9–1.10 & Logaritmické rovnice typu log\_a f(x)=b: podmínky. \\
37 & 1.6, 1.9–1.10 & Logaritmické rovnice se stejnými základy: porovnání argumentů. \\
38 & 1.6, 1.9–1.10 & Rovnice s použitím logaritmických vzorců. \\
39 & 1.6, 1.9–1.10 & Rovnice s více logaritmy a podmínkami pro argumenty. \\
40 & 1.6, 1.9–1.10 & Logaritmické rovnice s kvadratickým argumentem. \\
41 & T3 & \testr{Průběžný test 3 (15–20 min): Logaritmické rovnice. Zbytek hodiny: krátký rozbor a navazující procvičování.} \\
42 & 1.6, 1.9–1.10 & Logaritmické nerovnice: základ a směr nerovnosti. \\
\addlinespace[1.1mm]\multicolumn{2}{@{}l}{\textbf{Leden}} & \multicolumn{1}{r@{}}{\textbf{12 hodin}}\\
\midrule
43 & 1.6, 1.9–1.10 & Logaritmické nerovnice s více podmínkami. \\
44 & 1.6, 1.9–1.10 & Smíšené logaritmické rovnice a nerovnice. \\
45 & 1.6, 1.9–1.10 & Opakování exponenciálních a logaritmických rovnic. \\
46 & 1.6, 1.9–1.10 & Příklady s nepřípustnými kořeny a nezbytností zkoušky. \\
47 & 2.1–2.5 & Orientovaný úhel a jednotková kružnice. \\
48 & 2.1–2.5 & Stupňová a oblouková míra; převody a úhly lišící se o otočky. \\
49 & T4 & \testr{Průběžný test 4 (15–20 min): Exponenciální a logaritmické rovnice, základy goniometrie. Zbytek hodiny: krátký rozbor a navazující procvičování.} \\
50 & 2.1–2.5 & Sinus a kosinus obecného úhlu; znaménka v kvadrantech. \\
51 & 2.1–2.5 & Tangens a kotangens; definiční obor, nuly a znaménka. \\
52 & P2 & \testr{Pololetní písemná práce: Exponenciální a logaritmické rovnice, základy goniometrie.} \\
53 & 2.1–2.5 & Hodnoty goniometrických funkcí významných úhlů. \\
54 & 2.1–2.5 & Tabulka základních hodnot a určení úhlu z funkční hodnoty. \\
\addlinespace[1mm]\multicolumn{3}{@{}l}{\footnotesize\textbf{3. čtvrtletí} -- 2 průběžné testy + 1 souhrnná práce}\\[0.5mm]
\addlinespace[1.1mm]\multicolumn{2}{@{}l}{\textbf{Únor}} & \multicolumn{1}{r@{}}{\textbf{9 hodin}}\\
\midrule
55 & 2.1–2.5 & Graf funkce sinus: perioda, amplituda, obor hodnot. \\
56 & 2.1–2.5 & Graf funkce kosinus a transformace grafů. \\
57 & 2.1–2.5 & Grafy tangens a kotangens; asymptoty. \\
58 & 2.1–2.5 & Základní goniometrické vzorce a jejich význam. \\
59 & 2.1–2.5 & Vzorce pro součet a rozdíl úhlů; jednoduché úpravy. \\
60 & 2.1–2.5 & Vzorce pro dvojnásobný úhel a úpravy výrazů. \\
61 & T5 & \testr{Průběžný test 5 (15–20 min): Hodnoty a grafy goniometrických funkcí. Zbytek hodiny: krátký rozbor a navazující procvičování.} \\
62 & 2.1–2.5 & Goniometrické identity: ověření pomocí známých vztahů. \\
63 & 2.1–2.5 & Rovnice sin x=a, cos x=a: všechna řešení. \\
\addlinespace[1.1mm]\multicolumn{2}{@{}l}{\textbf{Březen}} & \multicolumn{1}{r@{}}{\textbf{12 hodin}}\\
\midrule
64 & 2.1–2.5 & Rovnice tg x=a a cotg x=a: periodické řešení. \\
65 & 2.1–2.5 & Goniometrické rovnice s argumentem kx+b. \\
66 & 2.1–2.5 & Goniometrické rovnice se substitucí a úpravou výrazů. \\
67 & 2.6–2.8 & Sinová věta: vztah mezi stranami a protilehlými úhly. \\
68 & 2.6–2.8 & Kosinová věta: výpočet strany obecného trojúhelníku. \\
69 & 2.6–2.8 & Kosinová věta: výpočet úhlu a kontrola výsledku. \\
70 & 2.6–2.8 & Volba mezi sinovou a kosinovou větou podle zadání. \\
71 & 2.6–2.8 & Obecný trojúhelník: dopočítání všech stran a úhlů. \\
72 & 2.6–2.8 & Obsah trojúhelníku a technické úlohy. \\
73 & 2.6–2.8 & Sinová a kosinová věta: smíšené příklady. \\
74 & 2.6–2.8 & Opakování trigonometrie před testem. \\
75 & T6 & \testr{Průběžný test 6 (15–20 min): Sinová a kosinová věta. Zbytek hodiny: krátký rozbor a navazující procvičování.} \\
\addlinespace[1.1mm]\multicolumn{2}{@{}l}{\textbf{Duben}} & \multicolumn{1}{r@{}}{\textbf{3 hodiny}}\\
\midrule
76 & 3.1–3.9 & Základní geometrické pojmy: body, přímky, úhly, rovnoběžnost. \\
77 & P3 & \testr{Čtvrtletní písemná práce: Goniometrie, trigonometrie a technické výpočty.} \\
78 & 3.1–3.9 & Trojúhelníky: úhly, těžnice, výšky, osy. \\
\addlinespace[1mm]\multicolumn{3}{@{}l}{\footnotesize\textbf{4. čtvrtletí} -- 2 průběžné testy + 1 souhrnná práce}\\[0.5mm]
79 & 3.1–3.9 & Kružnice a kruh: tečny, sečny, vzájemná poloha. \\
80 & 3.1–3.9 & Shodnost trojúhelníků a věty o shodnosti. \\
81 & 3.1–3.9 & Podobnost trojúhelníků a dělení úsečky v poměru. \\
82 & 3.1–3.9 & Stejnolehlost: vlastnosti a jednoduché konstrukce. \\
83 & 3.1–3.9 & Pythagorova věta a její obrácení – procvičení. \\
84 & 3.1–3.9 & Euklidovy věty o výšce a odvěsnách. \\
\addlinespace[1.1mm]\multicolumn{2}{@{}l}{\textbf{Květen}} & \multicolumn{1}{r@{}}{\textbf{9 hodin}}\\
\midrule
85 & 3.1–3.9 & Výpočty v pravoúhlém trojúhelníku; vhodná metoda. \\
86 & 3.1–3.9 & Konstrukce trojúhelníku: náčrt a rozbor. \\
87 & 3.1–3.9 & Konstrukce trojúhelníku z kombinace stran a úhlů. \\
88 & T7 & \testr{Průběžný test 7 (15–20 min): Podobnost, Pythagorova a Euklidovy věty. Zbytek hodiny: krátký rozbor a navazující procvičování.} \\
89 & 3.1–3.9 & Konstrukce podle těžnice a výšky; rozbor podmínek. \\
90 & 3.1–3.9 & Množiny bodů dané vlastnosti: osa úsečky a úhlu. \\
91 & 3.1–3.9 & Konstrukční úlohy s kružnicemi a jejich průsečíky. \\
92 & 3.1–3.9 & Středový a obvodový úhel; vztah ke kružnici. \\
93 & 3.1–3.9 & Obvody a obsahy rovinných útvarů; jednotky. \\
\addlinespace[1.1mm]\multicolumn{2}{@{}l}{\textbf{Červen}} & \multicolumn{1}{r@{}}{\textbf{9 hodin}}\\
\midrule
94 & 3.1–3.9 & Složené planimetrické úlohy a opakování. \\
95 & T8 & \testr{Průběžný test 8 (15–20 min): Konstrukce a planimetrie. Zbytek hodiny: krátký rozbor a navazující procvičování.} \\
96 & 4.1–4.4 & Vzájemná poloha bodů, přímek a rovin v prostoru. \\
97 & 4.1–4.4 & Zobrazení těles a prostorová představivost; řezy. \\
98 & 4.1–4.4 & Povrch a objem hranolu a válce. \\
99 & 4.1–4.4 & Povrch a objem jehlanu, kužele a koule. \\
100 & 4.1–4.4 & Komolá tělesa a jednoduché technické úlohy. \\
101 & P4 & \testr{Pololetní písemná práce: Planimetrie, základy stereometrie; souhrn 2. ročníku.} \\
102 & 4.1–4.4 & Závěrečné opakování geometrie a funkcí. \\
\end{longtable}
\normalsize
\clearpage
\section*{3. ročník – komplexní čísla a analytická geometrie}
\textit{Dotace 102 hodin. Čtvrtletní celky a tři testy v každém čtvrtletí jsou součástí hodinové dotace.}
\textit{Historická orientace: 3.A a 3.D postupovaly stejně po hlavních tématech, ale měly odlišná čísla hodin. Plán nepřiřazuje konkrétní historická čísla plánovaným hodinám.}
\vspace{2mm}
\footnotesize
\begin{longtable}{@{}>{\RaggedLeft\arraybackslash}p{9mm}>{\RaggedRight\arraybackslash}p{15mm}>{\RaggedRight\arraybackslash}p{150mm}@{}}
\toprule \textbf{Hod.} & \textbf{TP} & \textbf{Náplň vyučovací hodiny} \\\midrule\endfirsthead
\toprule \textbf{Hod.} & \textbf{TP} & \textbf{Náplň vyučovací hodiny} \\\midrule\endhead
\bottomrule\endfoot
\addlinespace[1mm]\multicolumn{3}{@{}l}{\footnotesize\textbf{1. čtvrtletí} -- 2 průběžné testy + 1 souhrnná práce}\\[0.5mm]
\addlinespace[1.1mm]\multicolumn{2}{@{}l}{\textbf{Září}} & \multicolumn{1}{r@{}}{\textbf{9 hodin}}\\
\midrule
1 & 1.1–1.7 & Úvod do třetího ročníku, opakování druhého ročníku. \\
2 & 1.1–1.7 & Zavedení komplexního čísla, imaginární jednotka a algebraický tvar. \\
3 & 1.1–1.7 & Sčítání a odčítání komplexních čísel. \\
4 & 1.1–1.7 & Násobení komplexních čísel v algebraickém tvaru. \\
5 & 1.1–1.7 & Komplexně sdružené číslo; sdruženost a absolutní hodnota. \\
6 & 1.1–1.7 & Dělení komplexních čísel a usměrnění jmenovatele. \\
7 & 1.1–1.7 & Operace s komplexními čísly – procvičení. \\
8 & 1.1–1.7 & Gaussova rovina: obraz komplexního čísla, modul, argument. \\
9 & 1.1–1.7 & Převod algebraického tvaru na goniometrický. \\
\addlinespace[1.1mm]\multicolumn{2}{@{}l}{\textbf{Říjen}} & \multicolumn{1}{r@{}}{\textbf{12 hodin}}\\
\midrule
10 & T1 & \testr{Průběžný test 1 (15–20 min): Operace s komplexními čísly v algebraickém tvaru. Zbytek hodiny: krátký rozbor a navazující procvičování.} \\
11 & 1.1–1.7 & Goniometrický tvar komplexního čísla: argument a modul. \\
12 & 1.1–1.7 & Násobení a dělení čísel v goniometrickém tvaru. \\
13 & 1.1–1.7 & Moivreova věta: mocniny komplexních čísel. \\
14 & 1.1–1.7 & Odmocňování komplexních čísel – geometrický význam. \\
15 & 1.1–1.7 & Kvadratická rovnice v oboru komplexních čísel. \\
16 & 1.1–1.7 & Binomické rovnice: základní příklady. \\
17 & 1.1–1.7 & Komplexní čísla: převody tvarů, operace a kontrola. \\
18 & 1.1–1.7 & Opakování komplexních čísel před souhrnným testem. \\
19 & 2.1–2.2 & Analytická geometrie v rovině: kartézská soustava souřadnic. \\
20 & T2 & \testr{Průběžný test 2 (15–20 min): Komplexní čísla – všechny probírané tvary. Zbytek hodiny: krátký rozbor a navazující procvičování.} \\
21 & 2.1–2.2 & Vzdálenost dvou bodů v rovině; vzorec a Pythagorova věta. \\
\addlinespace[1.1mm]\multicolumn{2}{@{}l}{\textbf{Listopad}} & \multicolumn{1}{r@{}}{\textbf{7 hodin}}\\
\midrule
22 & 2.1–2.2 & Střed úsečky: souřadnice a slovní úlohy. \\
23 & 2.1–2.2 & Trojúhelník v souřadnicích: délky stran a těžnice. \\
24 & 2.1–2.2 & Vektor jako orientovaná úsečka a jeho souřadnice. \\
25 & 2.1–2.2 & Velikost vektoru a jednotkový vektor. \\
26 & 2.1–2.2 & Sčítání a odčítání vektorů; geometrický význam. \\
27 & P1 & \testr{Čtvrtletní písemná práce: Komplexní čísla a úvod analytické geometrie.} \\
28 & 2.2–2.3 & Násobení vektoru reálným číslem; rovnoběžnost vektorů. \\
\addlinespace[1mm]\multicolumn{3}{@{}l}{\footnotesize\textbf{2. čtvrtletí} -- 2 průběžné testy + 1 souhrnná práce}\\[0.5mm]
29 & 2.2–2.3 & Lineární kombinace vektorů a lineární závislost. \\
30 & 2.2–2.3 & Skalární součin: definice a výpočet ze souřadnic. \\
31 & 2.2–2.3 & Kolmost a odchylka vektorů pomocí skalárního součinu. \\
32 & 2.2–2.3 & Skalární součin: smíšené úlohy a ověřování. \\
33 & 2.2–2.3 & Vektorový součin vektorů – význam a výpočet. \\
\addlinespace[1.1mm]\multicolumn{2}{@{}l}{\textbf{Prosinec}} & \multicolumn{1}{r@{}}{\textbf{9 hodin}}\\
\midrule
34 & 2.2–2.3 & Geometrický význam vektorového součinu: kolmý vektor. \\
35 & 2.2–2.3 & Obsah rovnoběžníku a trojúhelníku z vektorového součinu. \\
36 & T3 & \testr{Průběžný test 3 (15–20 min): Skalární a vektorový součin. Zbytek hodiny: krátký rozbor a navazující procvičování.} \\
37 & 2.2–2.3 & Vektorové operace: postup výpočtu a kontrola. \\
38 & 2.2–2.3 & Skalární a vektorový součin: výběr správné operace. \\
39 & 2.2–2.3 & Opakování vektorů a součinů před testem. \\
40 & 2.2–2.3 & Samostatné procvičení analytických úloh s vektory. \\
41 & 2.4–2.8 & Parametrické vyjádření přímky: bod a směrový vektor. \\
42 & 2.4–2.8 & Rovnice přímky z dvojice bodů: nalezení směrového vektoru. \\
\addlinespace[1.1mm]\multicolumn{2}{@{}l}{\textbf{Leden}} & \multicolumn{1}{r@{}}{\textbf{12 hodin}}\\
\midrule
43 & 2.4–2.8 & Bod na přímce: dosazení do parametrického tvaru. \\
44 & 2.4–2.8 & Parametrický tvar přímky: graf a procvičování. \\
45 & T4 & \testr{Průběžný test 4 (15–20 min): Vektory a parametrické vyjádření přímky. Zbytek hodiny: krátký rozbor a navazující procvičování.} \\
46 & 2.4–2.8 & Obecná rovnice přímky: normálový vektor. \\
47 & 2.4–2.8 & Převod parametrického vyjádření na obecnou rovnici. \\
48 & 2.4–2.8 & Převod obecné rovnice na parametrické vyjádření. \\
49 & 2.4–2.8 & Obecná rovnice přímky: poloha bodu vůči přímce. \\
50 & 2.4–2.8 & Směrnicový tvar přímky a jeho geometrický význam. \\
51 & 2.4–2.8 & Převody mezi třemi rovnicemi přímky. \\
52 & 2.4–2.8 & Vzájemná poloha dvou přímek: rovnoběžnost, totožnost, různoběžnost. \\
53 & P2 & \testr{Pololetní písemná práce: Komplexní čísla, vektory a přímky v rovině.} \\
54 & 2.4–2.8 & Výpočet průsečíku přímek z jejich rovnic. \\
\addlinespace[1mm]\multicolumn{3}{@{}l}{\footnotesize\textbf{3. čtvrtletí} -- 2 průběžné testy + 1 souhrnná práce}\\[0.5mm]
\addlinespace[1.1mm]\multicolumn{2}{@{}l}{\textbf{Únor}} & \multicolumn{1}{r@{}}{\textbf{9 hodin}}\\
\midrule
55 & 2.4–2.8 & Kolmost přímek a normálové vektory. \\
56 & 2.4–2.8 & Odchylka dvou přímek ze směrových vektorů. \\
57 & 2.4–2.8 & Vzdálenost bodu od přímky: odvození a výpočet. \\
58 & 2.4–2.8 & Vzdálenost dvou rovnoběžných přímek. \\
59 & 2.4–2.8 & Přímky v rovině: slovní a technické úlohy. \\
60 & 2.4–2.8 & Analytická geometrie přímek: úlohy s parametrem. \\
61 & 2.4–2.8 & Opakování rovnic přímky a vzájemné polohy. \\
62 & 2.4–2.8 & Opakování odchylky a vzdálenosti v rovině. \\
63 & 3.1–3.4 & Analytická geometrie v prostoru: souřadnice bodu a vektoru. \\
\addlinespace[1.1mm]\multicolumn{2}{@{}l}{\textbf{Březen}} & \multicolumn{1}{r@{}}{\textbf{12 hodin}}\\
\midrule
64 & 3.1–3.4 & Vzdálenost bodů a velikost vektoru v prostoru. \\
65 & T5 & \testr{Průběžný test 5 (15–20 min): Rovnice přímek, odchylky a vzdálenosti. Zbytek hodiny: krátký rozbor a navazující procvičování.} \\
66 & 3.1–3.4 & Parametrické vyjádření přímky v prostoru. \\
67 & 3.1–3.4 & Vzájemná poloha přímek v prostoru: různé případy. \\
68 & 3.1–3.4 & Mimoběžky a jejich rozpoznání pomocí rovnic. \\
69 & 3.1–3.4 & Parametrické vyjádření roviny: bod a dva směrové vektory. \\
70 & 3.1–3.4 & Obecná rovnice roviny a normálový vektor. \\
71 & 3.1–3.4 & Vzájemná poloha přímky a roviny; průsečík. \\
72 & 3.1–3.4 & Vzájemná poloha dvou rovin; průsečnice. \\
73 & 3.1–3.4 & Vzdálenost bodu od roviny – vzorec a podmínky. \\
74 & 3.1–3.4 & Odchylka přímek v prostoru: směrové vektory. \\
75 & T6 & \testr{Průběžný test 6 (15–20 min): Analytická geometrie v prostoru. Zbytek hodiny: krátký rozbor a navazující procvičování.} \\
\addlinespace[1.1mm]\multicolumn{2}{@{}l}{\textbf{Duben}} & \multicolumn{1}{r@{}}{\textbf{3 hodiny}}\\
\midrule
76 & 3.1–3.4 & Odchylka přímky a roviny; odchylka rovin. \\
77 & 3.1–3.4 & Vzdálenosti v prostoru: procvičení typových úloh. \\
78 & P3 & \testr{Čtvrtletní písemná práce: Přímky a roviny v prostoru, vzdálenosti.} \\
\addlinespace[1mm]\multicolumn{3}{@{}l}{\footnotesize\textbf{4. čtvrtletí} -- 2 průběžné testy + 1 souhrnná práce}\\[0.5mm]
79 & 4.1–4.9 & Kuželosečky: geometrická definice a společné souvislosti. \\
80 & 4.1–4.9 & Kružnice: středová rovnice a základní vlastnosti. \\
81 & 4.1–4.9 & Kružnice: obecná rovnice a doplnění na čtverec. \\
82 & 4.1–4.9 & Kružnice a přímka: sečna, tečna, vnější přímka. \\
83 & 4.1–4.9 & Tečna ke kružnici v daném bodě; normála. \\
84 & 4.1–4.9 & Kružnice: úlohy s parametrem a polohou přímky. \\
\addlinespace[1.1mm]\multicolumn{2}{@{}l}{\textbf{Květen}} & \multicolumn{1}{r@{}}{\textbf{9 hodin}}\\
\midrule
85 & 4.1–4.9 & Elipsa: ohnisková definice a geometrická konstrukce. \\
86 & 4.1–4.9 & Středová rovnice elipsy; poloosy a výstřednost. \\
87 & 4.1–4.9 & Elipsa z rovnice: určení středu, poloos a ohnisek. \\
88 & T7 & \testr{Průběžný test 7 (15–20 min): Kružnice a elipsa. Zbytek hodiny: krátký rozbor a navazující procvičování.} \\
89 & 4.1–4.9 & Elipsa: sestavení rovnice ze zadaných údajů. \\
90 & 4.1–4.9 & Přímka a elipsa: počet průsečíků, dosazení. \\
91 & 4.1–4.9 & Kružnice a elipsa: souhrnné procvičení a porovnání. \\
92 & 4.1–4.9 & Parabola: definice pomocí ohniska a řídicí přímky. \\
93 & 4.1–4.9 & Parabola: vrcholový tvar rovnice a orientace osy. \\
\addlinespace[1.1mm]\multicolumn{2}{@{}l}{\textbf{Červen}} & \multicolumn{1}{r@{}}{\textbf{9 hodin}}\\
\midrule
94 & 4.1–4.9 & Parabola z rovnice: ohnisko, řídicí přímka a náčrt. \\
95 & 4.1–4.9 & Přímka a parabola: průsečíky a tečna. \\
96 & 4.1–4.9 & Hyperbola: ohnisková definice, vrcholy a asymptoty. \\
97 & 4.1–4.9 & Hyperbola: základní rovnice a geometrický náčrt. \\
98 & 4.1–4.9 & Hyperbola: určení středu a poloos z rovnice. \\
99 & 4.1–4.9 & Srovnání elipsy, paraboly a hyperboly; rozlišovací znaky. \\
100 & T8 & \testr{Průběžný test 8 (15–20 min): Kuželosečky: parabola a hyperbola. Zbytek hodiny: krátký rozbor a navazující procvičování.} \\
101 & 4.1–4.9 & Kuželosečky: smíšené úlohy a opakování. \\
102 & P4 & \testr{Pololetní písemná práce: Kuželosečky a opakování analytické geometrie.} \\
\end{longtable}
\normalsize
\clearpage
\section*{4. ročník – posloupnosti, pravděpodobnost, statistika a opakování}
\textit{Dotace 112 hodin. Čtvrtletní celky a tři testy v každém čtvrtletí jsou součástí hodinové dotace.}
\textit{Obsah 4. ročníku není doložen historickými zápisy; vychází z TP 2020 a dřívější osnovy.}
\vspace{2mm}
\footnotesize
\begin{longtable}{@{}>{\RaggedLeft\arraybackslash}p{9mm}>{\RaggedRight\arraybackslash}p{15mm}>{\RaggedRight\arraybackslash}p{150mm}@{}}
\toprule \textbf{Hod.} & \textbf{TP} & \textbf{Náplň vyučovací hodiny} \\\midrule\endfirsthead
\toprule \textbf{Hod.} & \textbf{TP} & \textbf{Náplň vyučovací hodiny} \\\midrule\endhead
\bottomrule\endfoot
\addlinespace[1mm]\multicolumn{3}{@{}l}{\footnotesize\textbf{1. čtvrtletí} -- 2 průběžné testy + 1 souhrnná práce}\\[0.5mm]
\addlinespace[1.1mm]\multicolumn{2}{@{}l}{\textbf{Září}} & \multicolumn{1}{r@{}}{\textbf{12 hodin}}\\
\midrule
1 & 1.1 & Posloupnost jako funkce na přirozených číslech; členy a značení. \\
2 & 1.1 & Výpis členů a vzorec pro n-tý člen; hledání předpisu. \\
3 & 1.1 & Graf posloupnosti a jeho diskrétní charakter. \\
4 & 1.1 & Rekurentní zadání posloupnosti; výpočet členů. \\
5 & 1.1 & Explicitní a rekurentní zadání – srovnání a úlohy. \\
6 & 1.1 & Vlastnosti posloupnosti: monotónnost a omezenost. \\
7 & 1.2 & Aritmetická posloupnost; diference a vztah pro n-tý člen. \\
8 & 1.2 & Slovní úlohy a procvičení aritmetické posloupnosti. \\
9 & 1.3 & Geometrická posloupnost; kvocient a vztah pro n-tý člen. \\
10 & 1.3 & Součet prvních n členů geometrické posloupnosti. \\
11 & 1.3 & Procvičení geometrické posloupnosti; rozdíl oproti aritmetické. \\
12 & T1 & \testr{Průběžný test 1 (15–20 min): Posloupnosti a jejich zadávání. Zbytek hodiny: krátký rozbor a navazující procvičování.} \\
\addlinespace[1.1mm]\multicolumn{2}{@{}l}{\textbf{Říjen}} & \multicolumn{1}{r@{}}{\textbf{16 hodin}}\\
\midrule
13 & 1.4 & Užití posloupností v praxi – pravidelný přírůstek a rozdíl. \\
14 & 1.4 & Geometrický růst a pokles v technice a přírodě. \\
15 & 1.4 & Úlohy na růst počtu výrobků, výkonu či spotřeby. \\
16 & 1.4 & Srovnávací slovní úlohy na oba typy posloupností. \\
17 & 1.5 & Základní finanční pojmy: kapitál, úroková míra, období. \\
18 & 1.5 & Jednoduché úročení; výpočet úroku a konečné částky. \\
19 & 1.5 & Složené úrokování: odvození vzorce z geometrické posloupnosti. \\
20 & 1.5 & Složené úrokování: výpočet počátečního vkladu. \\
21 & 1.5 & Složené úrokování: časové období a úroková míra. \\
22 & 1.5 & Různé délky úrokovacího období; porovnání nabídek. \\
23 & 1.5 & Pravidelné vklady a spoření – jednoduchý model. \\
24 & 1.5 & Úvěry, splácení a praktická finanční gramotnost. \\
25 & T2 & \testr{Průběžný test 2 (15–20 min): Aritmetická, geometrická posloupnost a finance. Zbytek hodiny: krátký rozbor a navazující procvičování.} \\
26 & 1.4–1.5 & Smíšené aplikační úlohy na posloupnosti a úročení. \\
27 & 1.1–1.5 & Opakování před 1. písemnou prací. \\
28 & P1 & \testr{Čtvrtletní písemná práce: Posloupnosti a finanční matematika.} \\
\addlinespace[1mm]\multicolumn{3}{@{}l}{\footnotesize\textbf{2. čtvrtletí} -- 2 průběžné testy + 1 souhrnná práce}\\[0.5mm]
\addlinespace[1.1mm]\multicolumn{2}{@{}l}{\textbf{Listopad}} & \multicolumn{1}{r@{}}{\textbf{16 hodin}}\\
\midrule
29 & 2.1 & Úvod do kombinatoriky: systematický výpis, význam pořadí a opakování. \\
30 & 2.1 & Pravidlo součtu a součinu; strom možností, více kroků výběru. \\
31 & 2.2 & Faktoriál: definice, zkracování a úpravy jednoduchých výrazů. \\
32 & 2.1 & Permutace bez opakování; odvození počtu uspořádání. \\
33 & 2.1 & Permutace s opakováním; počet různých pořadí stejných prvků. \\
34 & 2.1 & Variace s opakováním; PINy, kódy a posloupnosti znaků. \\
35 & 2.1 & Kombinace bez opakování; význam neuspořádaného výběru. \\
36 & 2.2 & Kombinační číslo a jeho výpočet pomocí faktoriálu. \\
37 & 2.2 & Vlastnosti kombinačních čísel; symetrie a Pascalův vztah. \\
38 & 2.3 & Pascalův trojúhelník a jeho souvislost s kombinačními čísly. \\
39 & 2.3 & Binomická věta: rozvoj (a+b)\textasciicircum{}n pro malé exponenty. \\
40 & T3 & \testr{Průběžný test 3 (15–20 min): Variace, permutace a kombinace. Zbytek hodiny: krátký rozbor a navazující procvičování.} \\
41 & 2.1–2.3 & Smíšené kombinatorické úlohy: rozpoznání typu a řešení. \\
42 & 2.4 & Elementární jevy a výpis výsledků; hod mincí, kostkou. \\
43 & 2.4 & Zobrazení výsledkového prostoru pomocí tabulky či stromu. \\
44 & 2.5 & Náhodný jev: jistý, nemožný, opačný; množinové vyjádření. \\
\addlinespace[1.1mm]\multicolumn{2}{@{}l}{\textbf{Prosinec}} & \multicolumn{1}{r@{}}{\textbf{12 hodin}}\\
\midrule
45 & 2.5 & Klasická pravděpodobnost, stejně pravděpodobné výsledky. \\
46 & 2.5 & Pravděpodobnost v úlohách s kostkou, mincí a kartami. \\
47 & 2.5 & Pravděpodobnost s použitím kombinatorického počítání. \\
48 & 2.6 & Sjednocení a průnik jevů; výklad pomocí diagramů. \\
49 & 2.6 & Vzorec pro pravděpodobnost sjednocení dvou jevů. \\
50 & 2.6 & Pravděpodobnost opačného jevu; aplikační úlohy. \\
51 & 2.7 & Procvičování nezávislých jevů v jednoduchých pokusech. \\
52 & 2.8 & Statistický soubor, statistická jednotka a znak. \\
53 & T4 & \testr{Průběžný test 4 (15–20 min): Kombinatorika a základní pravděpodobnost. Zbytek hodiny: krátký rozbor a navazující procvičování.} \\
54 & 2.9 & Absolutní a relativní četnost; tabulka četností. \\
55 & 2.9 & Četnostní grafy a interpretace dat – praktické příklady. \\
56 & P2 & \testr{Pololetní písemná práce: Kombinatorika a základní pravděpodobnost.} \\
\addlinespace[1mm]\multicolumn{3}{@{}l}{\footnotesize\textbf{3. čtvrtletí} -- 2 průběžné testy + 1 souhrnná práce}\\[0.5mm]
\addlinespace[1.1mm]\multicolumn{2}{@{}l}{\textbf{Leden}} & \multicolumn{1}{r@{}}{\textbf{16 hodin}}\\
\midrule
57 & 2.10 & Charakteristiky polohy; aritmetický průměr. \\
58 & 2.10 & Vážený aritmetický průměr a práce s četnostmi. \\
59 & 2.10 & Modus a medián; význam odlehlých hodnot. \\
60 & 2.10 & Porovnání průměru, mediánu a modu na datech. \\
61 & 2.10 & Rozptyl – princip a výpočet pro datový soubor. \\
62 & 2.10 & Směrodatná odchylka a její interpretace. \\
63 & 2.10 & Výpočet rozptylu a odchylky pomocí kalkulačky. \\
64 & 2.10 & Charakteristiky variability z tabulky četností. \\
65 & 2.10 & Technická statistika: měření, tolerance a rozptýlení hodnot. \\
66 & 2.10 & Zpracování jednoduchého statistického šetření. \\
67 & 2.10 & Interpretace a kritické čtení statistických souhrnů. \\
68 & 2.5–2.10 & Procvičování pravděpodobnosti a statistiky. \\
69 & 2.1–2.10 & Opakování kombinatoriky, pravděpodobnosti a statistiky. \\
70 & T5 & \testr{Průběžný test 5 (15–20 min): Pravděpodobnost jevů a jejich vztahy. Zbytek hodiny: krátký rozbor a navazující procvičování.} \\
71 & 2.1–2.10 & Rozbor 2. písemné práce a odstranění nejčastějších chyb. \\
72 & 3.1 & Opakování číselných oborů, procent a poměru. \\
\addlinespace[1.1mm]\multicolumn{2}{@{}l}{\textbf{Únor}} & \multicolumn{1}{r@{}}{\textbf{12 hodin}}\\
\midrule
73 & 3.1 & Mocniny, odmocniny a úpravy výrazů. \\
74 & 3.1 & Mnohočleny, rozklad na součin, lomené výrazy. \\
75 & 3.1 & Lineární rovnice a nerovnice. \\
76 & 3.1 & Kvadratická rovnice a diskriminant. \\
77 & 3.1 & Kvadratická funkce, grafy a vrchol paraboly. \\
78 & 3.1 & Kvadratické nerovnice a odmocninové rovnice. \\
79 & 3.1 & Funkce – definiční obor, obor hodnot, graf, vlastnosti. \\
80 & T6 & \testr{Průběžný test 6 (15–20 min): Deskriptivní statistika a charakteristiky. Zbytek hodiny: krátký rozbor a navazující procvičování.} \\
81 & 3.1 & Exponenciální a logaritmické funkce. \\
82 & 3.1 & Exponenciální a logaritmické rovnice a nerovnice. \\
83 & 3.1 & Smíšený testový blok: algebra a funkce – rozbor postupů. \\
84 & P3 & \testr{Čtvrtletní písemná práce: Pravděpodobnost a statistika.} \\
\addlinespace[1mm]\multicolumn{3}{@{}l}{\footnotesize\textbf{4. čtvrtletí} -- 2 průběžné testy + 1 souhrnná práce}\\[0.5mm]
\addlinespace[1.1mm]\multicolumn{2}{@{}l}{\textbf{Březen}} & \multicolumn{1}{r@{}}{\textbf{16 hodin}}\\
\midrule
85 & 3.1 & Orientovaný úhel; převody radiánů a stupňů. \\
86 & 3.1 & Goniometrické funkce, grafy a vlastnosti. \\
87 & 3.1 & Goniometrické vzorce a úpravy výrazů. \\
88 & 3.1 & Goniometrické rovnice a interpretace řešení. \\
89 & 3.1 & Sinová a kosinová věta – obecný trojúhelník. \\
90 & 3.1 & Planimetrie: trojúhelníky, shodnost, podobnost. \\
91 & 3.1 & Planimetrie: kružnice, úhly, obsahy a obvody. \\
92 & 3.1 & Stereometrie: povrchy a objemy těles. \\
93 & 3.1 & Analytická geometrie: body a vektory v rovině. \\
94 & 3.1 & Analytická geometrie: rovnice přímky, odchylky a vzdálenosti. \\
95 & 3.1 & Analytická geometrie: přímky a roviny v prostoru. \\
96 & 3.1 & Kuželosečky: kružnice, elipsa, hyperbola a parabola. \\
97 & T7 & \testr{Průběžný test 7 (15–20 min): Opakování funkcí a rovnic. Zbytek hodiny: krátký rozbor a navazující procvičování.} \\
98 & 3.1 & Procvičování geometrických a trigonometrických úloh. \\
99 & 3.1 & Opakování před 3. písemnou prací. \\
100 & 3.1 & Rozbor 3. písemné práce a individuální oprava typických chyb. \\
\addlinespace[1.1mm]\multicolumn{2}{@{}l}{\textbf{Duben}} & \multicolumn{1}{r@{}}{\textbf{12 hodin}}\\
\midrule
101 & 3.1 & Komplexní čísla – algebraický a goniometrický tvar. \\
102 & 3.1 & Komplexní čísla – operace, mocniny a kořeny. \\
103 & 3.1 & Posloupnosti, finanční matematika – souhrnné úlohy. \\
104 & 3.1 & Pravděpodobnost a statistika – smíšené úlohy. \\
105 & 3.1 & Maturitní úlohy: algebraické výrazy a rovnice. \\
106 & T8 & \testr{Průběžný test 8 (15–20 min): Opakování geometrie a kombinatoriky. Zbytek hodiny: krátký rozbor a navazující procvičování.} \\
107 & 3.1 & Maturitní úlohy: funkce a grafy. \\
108 & 3.1 & Maturitní úlohy: planimetrie a stereometrie. \\
109 & 3.1 & Maturitní úlohy: analytická geometrie a kuželosečky. \\
110 & 3.1 & Cvičný souhrnný didaktický test či blok úloh. \\
111 & P4 & \testr{Pololetní písemná práce: Souhrnné opakování matematiky před maturitou.} \\
112 & 3.1 & Rozbor testu, závěrečné upevnění a individuální doporučení. \\
\end{longtable}
\normalsize
\clearpage
\section*{Souhrn písemného ověřování}
Přehled pro rychlé vyhledání. Krátké testy lze zařadit do první poloviny hodiny; souhrnné práce jsou plánovány jako celohodinové.
\footnotesize
\begin{longtable}{@{}p{17mm}p{12mm}p{27mm}p{26mm}p{72mm}@{}}
\toprule \textbf{Ročník} & \textbf{Čtvrt.} & \textbf{Čísla hodin} & \textbf{Druh} & \textbf{Rozsah}\\\midrule\endfirsthead
\toprule \textbf{Ročník} & \textbf{Čtvrt.} & \textbf{Čísla hodin} & \textbf{Druh} & \textbf{Rozsah}\\\midrule\endhead
\bottomrule\endfoot
1. & 1. & 10, 21, 32 & 2 průběžné + 1 souhrnná & Výroková logika a číselné obory; Množiny, intervaly a množinové operace; Výroková logika, množiny, goniometrie a mocniny \\
1. & 2. & 40, 57, 69 & 2 průběžné + 1 souhrnná & Mocniny a odmocniny; Algebraické výrazy a rozklady; Algebraické výrazy a základy funkcí \\
1. & 3. & 84, 98, 110 & 2 průběžné + 1 souhrnná & Lineární rovnice a nerovnice; Soustavy a rovnice s absolutní hodnotou; Lineární funkce a rovnice, úvod kvadratických funkcí \\
1. & 4. & 121, 130, 135 & 2 průběžné + 1 souhrnná & Kvadratické rovnice a grafy; Kvadratické nerovnice a odmocninové rovnice; Souhrnné učivo 1. ročníku \\
2. & 1. & 9, 20, 29 & 2 průběžné + 1 souhrnná & Funkce a inverzní funkce; Exponenciální funkce a grafy; Funkce a exponenciální rovnice \\
2. & 2. & 41, 49, 52 & 2 průběžné + 1 souhrnná & Logaritmické rovnice; Exponenciální a logaritmické rovnice, základy goniometrie; Exponenciální a logaritmické rovnice, základy goniometrie \\
2. & 3. & 61, 75, 77 & 2 průběžné + 1 souhrnná & Hodnoty a grafy goniometrických funkcí; Sinová a kosinová věta; Goniometrie, trigonometrie a technické výpočty \\
2. & 4. & 88, 95, 101 & 2 průběžné + 1 souhrnná & Podobnost, Pythagorova a Euklidovy věty; Konstrukce a planimetrie; Planimetrie, základy stereometrie; souhrn 2. ročníku \\
3. & 1. & 10, 20, 27 & 2 průběžné + 1 souhrnná & Operace s komplexními čísly v algebraickém tvaru; Komplexní čísla – všechny probírané tvary; Komplexní čísla a úvod analytické geometrie \\
3. & 2. & 36, 45, 53 & 2 průběžné + 1 souhrnná & Skalární a vektorový součin; Vektory a parametrické vyjádření přímky; Komplexní čísla, vektory a přímky v rovině \\
3. & 3. & 65, 75, 78 & 2 průběžné + 1 souhrnná & Rovnice přímek, odchylky a vzdálenosti; Analytická geometrie v prostoru; Přímky a roviny v prostoru, vzdálenosti \\
3. & 4. & 88, 100, 102 & 2 průběžné + 1 souhrnná & Kružnice a elipsa; Kuželosečky: parabola a hyperbola; Kuželosečky a opakování analytické geometrie \\
4. & 1. & 12, 25, 28 & 2 průběžné + 1 souhrnná & Posloupnosti a jejich zadávání; Aritmetická, geometrická posloupnost a finance; Posloupnosti a finanční matematika \\
4. & 2. & 40, 53, 56 & 2 průběžné + 1 souhrnná & Variace, permutace a kombinace; Kombinatorika a základní pravděpodobnost; Kombinatorika a základní pravděpodobnost \\
4. & 3. & 70, 80, 84 & 2 průběžné + 1 souhrnná & Pravděpodobnost jevů a jejich vztahy; Deskriptivní statistika a charakteristiky; Pravděpodobnost a statistika \\
4. & 4. & 97, 106, 111 & 2 průběžné + 1 souhrnná & Opakování funkcí a rovnic; Opakování geometrie a kombinatoriky; Souhrnné opakování matematiky před maturitou \\
\end{longtable}
\normalsize
\textbf{Pracovní povaha plánu:} Náplně jednotlivých hodin jsou novým didaktickým návrhem učitele s využitím podkladů. Nejsou přepisem záznamů třídních knih a netvrdí, že každá plánovaná hodina byla v minulosti skutečně odučena.
\end{document}
